\documentclass[11pt,a4paper,titlepage,oneside]{report}

%%% RELACIÓN DE VARIABLES A PERSONALIZAR %%%
% \def\lingua{gal}
\def\lingua{esp} % descomenta esta liña se redactarás a memoria en español
%\def\lingua{eng} % descomenta esta liña se redactarás a memoria en inglés
\def\nome{David Rodríguez Bacelar}                             % substitúe aquí o teu nome
\def\nomedirectorA{Patricia Martín Rodilla}              % substitúe aquí o nome de quen dirixe
%\def\nomedirectorB{Outro Nome Completo}             % duplica esta liña máis veces se o precisas, cambiando
% a letra final (A, B, C, D...): úsanse na portada.tex
\def\titulo{Título completo do traballo fin de grao} % substitúe aquí o título do teu TFG
%\def\titulacion{gced}                               % descomenta esta liña e comenta a seguinte se es estudante do GCED
\def\titulacion{gei}
%\def\mencion{COMPUTACIÓN}                           % descomenta a mención que che corresponda se es estudante do GEI
\def\mencion{ENXEÑARÍA DO SOFTWARE}
%\def\mencion{ENXEÑARÍA DE COMPUTADORES}
%\def\mencion{SISTEMAS DE INFORMACIÓN}
%\def\mencion{TECNOLOXÍAS DA INFORMACIÓN}

%\def\renomearcadros{si} % descomenta esta liña se redactas a memoria en español e prefires que
% os "cuadros" e o "índice de cuadros" se renomeen
% a "tablas" e "índice de tablas" respectivamente

\usepackage{estilo_tfg}

% Lista de paquetes potencialmente interesantes (uso baixo demanda)

% \usepackage{alltt}       % proporciona o entorno alltt, semellante a verbatim pero que respecta comandos
% \usepackage{enumitem}    % permite personalizar os entornos de lista
% \usepackage{eurofont}    % proporciona o comando \euro
% \usepackage{float}       % permite máis opcións para controlar obxectos flotantes (táboas, figuras)
% \usepackage{hhline}      % permite personalizar as liñas horizontais en arrays e táboas
\usepackage{longtable}   % permite construir táboas que ocupan máis dunha páxina
% \usepackage{lscape}      % permite colocar partes do documento en orientación apaisada
% \usepackage{moreverb}    % permite personalizar o entorno verbatim
% \usepackage{multirow}    % permite crear celdas que ocupan varias filas da mesma táboa
% \usepackage{pdfpages}    % permite insertar ficheiros en PDF no documento
% \usepackage{rotating}    % permite diferentes tipos de rotacións para figuras e táboas
% \usepackage{subcaption}  % permite a inclusión de varias subfiguras nunha figura
% \usepackage{tabu}        % permite táboas flexibles
% \usepackage{tabularx}    % permite táboas con columnas de anchura determinada

%%%%%%%%%%%%%%%%%%%%%%%%%%%%%%%%%%%%%%%%%%%%%%%%%%%%%%%%%%%%%%%%%%%%%%%%%%%%%%%%
% Corpo                                                                        %
%%%%%%%%%%%%%%%%%%%%%%%%%%%%%%%%%%%%%%%%%%%%%%%%%%%%%%%%%%%%%%%%%%%%%%%%%%%%%%%%

\begin{document}

%%%%%%%%%%%%%%%%%%%%%%%%%%%%%%%%%%%%%%%%
% Preliminares do documento            %
%%%%%%%%%%%%%%%%%%%%%%%%%%%%%%%%%%%%%%%%

\begin{titlepage}
  
  \hspace*{128pt}
  \textcolor{udcpink}{{\fontencoding{T1}\fontfamily{phv}\selectfont Facultade de Informática}}\\[-32pt]

  \begin{center}
    \includegraphics[scale=0.3]{imaxes/udc}\\[25pt]

    {\large TRABALLO FIN DE GRAO \\
            \nometitulacion \\
            \nomemencion } \\[10pt]

    \carimbo \\[25pt]

    \begin{huge}
      \begin{spacing}{1.3}
        \bfseries \titulo
      \end{spacing}
    \end{huge}
  \end{center}
  
  \vfill
  
  \begin{flushright}
    {\large
    \begin{tabular}{ll}
      {\bf Estudante:} & \nome \\
      {\bf Dirección:} & \nomedirectorA \\
%                      & \nomedirectorB \\ % duplica esta liña máis veces se o precisas, cambiando
                                           % a letra final (A, B, C, D...); define eses nomes no memoria_tfg.tex
    \end{tabular}}
  \end{flushright}
  \rightline{A Coruña, \datasimple.}
\end{titlepage}

\dedicatoria{Dedicatoria} % escribe neste comando o teu texto de dedicatoria
\paxinaenbranco\
\begin{agradecementos}
  \blindtext\                % substitúe este comando polo teu texto de agradecementos
\end{agradecementos}
%%%%%%%%%%%%%%%%%%%%%%%%%%%%%%%%%%%%%%%%%%%%%%%%%%%%%%%%%%%%%%%%%%%%%%%%%%%%%%%%

\pagestyle{empty}
\begin{abstract}
  Con el objetivo de proporcionar una herramienta útil basada en el perfilado automático de autores, a lo largo de este trabajo
  se profundizará más en este campo en el contexto de las redes sociales, analizando cuáles son los mejores algoritmos,
  las técnicas más utilizadas y las características más comúnes que se suelen emplear.
  Asimismo, se desarrollará, bajo la metodología SCRUM, una aplicación web en la que se permita a los usuarios perfilar los autores de una
  colección de textos propia —haciendo uso de algoritmos seleccionados en la fase previa de investigación— y mostrando los resultado del perfilado en un \textit{dashboard} intuitivo y accesible.
  Finalmente, para mostrar un caso de uso real de la aplicación se empleará la colección de referencia sobre el movimiento \textit{Black Lives Matter} (\#BLM),
  puesta a disposición para este TFG por los tutores del mismo, con el objetivo de analizar el perfil de los usuarios que participaron en los debates en las redes sociales sobre este movimiento
  de gran relevancia social.

  \vspace*{25pt}
  \begin{segundoresumo}
    To provide a useful tool based on automatic author profiling, throughout this work, will research deeper into this field in the context of social networks, analyzing the best algorithms,
    the most commonly used techniques, the most commonly used techniques, and the typical features that are often employed. Furthermore, under the SCRUM methodology, we will
    develop a web application, allowing users
    to profile authors of their own text datasets employing algorithms selected in the previous research phase— and displaying the profiling results on an intuitive and accessible dashboard.
    Finally, to illustrate a real use case of the application, the reference collection related to the "Black Lives Matter" movement (\#BLM) will be utilized. This collection has been made available for
    this Bachelor's Thesis by its supervisors, with the aim of analyzing the profiles of users who participated in social media discussions about this socially significant movement.
  \end{segundoresumo}
  \vspace*{25pt}
  \newpage
\begin{multicols}{2}
\begin{description}
\item [\palabraschaveprincipal:] \mbox{} \\[-20pt]
  \begin{itemize}
    \item \textit{Perfilado de autor}
    \item \textit{Redes sociales}
    \item \textit{Procesado de lenguaje natural}
    \item \textit{Aplicación web}
    \item \textit{Aprendizaje automático}
  \end{itemize}
\end{description}

\begin{description}
\item [\palabraschavesecundaria:] \mbox{} \\[-20pt]
  \begin{itemize}
    \item \textit{Author profiling}
    \item \textit{Social networks}
    \item \textit{Natural language processing}
    \item \textit{Web application}
    \item \textit{Machine learning}
  \end{itemize}
\end{description}
\end{multicols}

\end{abstract}
\pagestyle{fancy}

%%%%%%%%%%%%%%%%%%%%%%%%%%%%%%%%%%%%%%%%%%%%%%%%%%%%%%%%%%%%%%%%%%%%%%%%%%%%%%%%


\pagenumbering{roman}
\setcounter{page}{1}
\bstctlcite{IEEEexample:BSTcontrol}

\tableofcontents
\listoffigures
\listoftables
\clearpage

\pagenumbering{arabic}
\setcounter{page}{1}

%%%%%%%%%%%%%%%%%%%%%%%%%%%%%%%%%%%%%%%%
% Capítulos                            %
%%%%%%%%%%%%%%%%%%%%%%%%%%%%%%%%%%%%%%%%

\chapter{Introdución}
\label{chap:introducion}

\lettrine{P}{rimer} capítulo da memoria, onde xeralmente se exporán as
liñas mestras do traballo, os obxectivos, etc. Incluimos un par de
exemplos de citas~\cite{ErlangBook,ErlangWebBook} e de referencias
internas (sección \ref{sec:mostra}, páxina \pageref{sec:mostra}).

\Blindtext

\section{Sección de mostra}
\label{sec:mostra}

\Blindtext

\subsection{Subsección de mostra}

\Blindtext

\chapter{Contido demostrativo}
\label{chap:demo}

\lettrine{E}{ntre} a introdución e as conclusións, o documento conterá
tantos capítulos como sexa preciso, sempre con coidado de non rebasar
o límite de 80 páxinas fixado polo regulamento de TFGs.

Empregaremos éste de xeito demostrativo, para ilustrar o uso de elementos
habituais que poidan ser de utilidade\footnote{Por exemplo, isto é unha nota a
  pé de páxina.}.

\section{Inclusión de imaxes}

Se precisamos imaxes no noso documento, incluirémolas do xeito que se indica na
figura~\ref{fig:exemplo} (páxina~\pageref{fig:exemplo}). Se o facemos así,
\LaTeX\ ubicará cada imaxe no mellor lugar posible, lugar que pode variar a
medida que o documento vaia crecendo coa inclusión de máis texto e outros
elementos (máis imaxes, táboas, etc.).

\begin{figure}[hp!]
  \centering
  \includegraphics[width=0.75\textwidth]{imaxes/udc.png}
  \caption{Pé de imaxe descritivo}
  \label{fig:exemplo}
\end{figure}

Recoméndase almacenar os ficheiros gráficos no directorio \texttt{imaxes}.

\subsection{Inclusión de varias sub-imaxes}

Se precisamos inserir imaxes relacionadas, pode ser apropiado incluílas como
sub-figuras, do xeito que se pode apreciar na
figura~\ref{fig:exemplo-subfiguras} coas imaxes~\ref{fig:subfigura-rotada}
e~\ref{fig:subfigura-deformada}. Como se pode ver nos exemplos desta sección,
sempre é recomendable referirse ás imaxes pola súa referencia, xa que dese
xeito non dependemos de onde queden ubicados os elementos en cuestión.

\begin{figure}[hp!]
  \centering
  \begin{subfigure}[c]{0.3\textwidth}
    \includegraphics[angle=45,width=\textwidth]{imaxes/udc.png}
    \caption{Pé de subimaxe rotada}
    \label{fig:subfigura-rotada}
  \end{subfigure}
  \hspace{0.1\textwidth}
  \begin{subfigure}[c]{0.3\textwidth}
    \includegraphics[width=\textwidth,height=3cm]{imaxes/udc.png}
    \caption{Pé de subimaxe deformada}
    \label{fig:subfigura-deformada}
  \end{subfigure}
  \caption{Pé de imaxe xeral}
  \label{fig:exemplo-subfiguras}
\end{figure}

\section{Inclusión de táboas}

Rendimiento de los diferentes algoritmos encontrados con respecto al \textit{profiling} de edad y sexo.

\begin{table}[hp!]
  \centering
  \rowcolors{2}{white}{udcgray!25}
  \begin{tabular}{c|c}
    \rowcolor{udcpink!25}
    \textbf{Título de columna} & \textbf{Outro título de columna} \\\hline
    \textit{Título de fila}    & Contido de celda                 \\
    \textit{Título de fila}    & Contido de celda                 \\
    \textit{Título de fila}    & Contido de celda                 \\
    \textit{Título de fila}    & Contido de celda                 \\
    \textit{Título de fila}    & Contido de celda                 \\
    \textit{Título de fila}    & Contido de celda                 \\
  \end{tabular}
  \caption{Pé de táboa descritivo}
  \label{tab:exemplo}
\end{table}

Para táboas longas que ocupan varias páxinas, como é o caso da \ref{tab:longa},
recoméndase o uso do paquete \texttt{lontable}, descomentando a liña
correspondente no ficheiro raíz do proxecto (\verb+memoria_tfg.tex+).

\rowcolors{2}{white}{udcgray!25}
\begin{longtable}{l|r|c}
  \caption{Pé descritivo dunha táboa longa}
  \label{tab:longa} \\

  \rowcolor{udcpink!25}
  \textbf{Primeira columna} & \textbf{Segunda columna} & \textbf{Terceira columna} \\\hline
  \endfirsthead

  \multicolumn{3}{c}{\tablename\ \thetable{} -- {\small \textit{(vén da páxina anterior)}}} \\
  \rowcolor{udcpink!25}
  \textbf{Primeira columna} & \textbf{Segunda columna} & \textbf{Terceira columna} \\\hline
  \endhead

  \multicolumn{3}{c}{\dotfill{\small \textit{(continúa na páxina seguinte)}}\dotfill} \\
  \endfoot

  \endlastfoot

  Texto de exemplo & abcdef ghjijklmn & 123.456778 \\
  Texto de exemplo & abcdef ghjijklmn & 123.456778 \\
  Texto de exemplo & abcdef ghjijklmn & 123.456778 \\
  Texto de exemplo & abcdef ghjijklmn & 123.456778 \\
  Texto de exemplo & abcdef ghjijklmn & 123.456778 \\
  Texto de exemplo & abcdef ghjijklmn & 123.456778 \\
  Texto de exemplo & abcdef ghjijklmn & 123.456778 \\
  Texto de exemplo & abcdef ghjijklmn & 123.456778 \\
  Texto de exemplo & abcdef ghjijklmn & 123.456778 \\
  Texto de exemplo & abcdef ghjijklmn & 123.456778 \\
  Texto de exemplo & abcdef ghjijklmn & 123.456778 \\
  Texto de exemplo & abcdef ghjijklmn & 123.456778 \\
  Texto de exemplo & abcdef ghjijklmn & 123.456778 \\
  Texto de exemplo & abcdef ghjijklmn & 123.456778 \\
  Texto de exemplo & abcdef ghjijklmn & 123.456778 \\
  Texto de exemplo & abcdef ghjijklmn & 123.456778 \\
  Texto de exemplo & abcdef ghjijklmn & 123.456778 \\
  Texto de exemplo & abcdef ghjijklmn & 123.456778 \\
  Texto de exemplo & abcdef ghjijklmn & 123.456778 \\
  Texto de exemplo & abcdef ghjijklmn & 123.456778 \\
  Texto de exemplo & abcdef ghjijklmn & 123.456778 \\
  Texto de exemplo & abcdef ghjijklmn & 123.456778 \\
  Texto de exemplo & abcdef ghjijklmn & 123.456778 \\
  Texto de exemplo & abcdef ghjijklmn & 123.456778 \\
  Texto de exemplo & abcdef ghjijklmn & 123.456778 \\
  Texto de exemplo & abcdef ghjijklmn & 123.456778 \\
  Texto de exemplo & abcdef ghjijklmn & 123.456778 \\
  Texto de exemplo & abcdef ghjijklmn & 123.456778 \\
  Texto de exemplo & abcdef ghjijklmn & 123.456778 \\
  Texto de exemplo & abcdef ghjijklmn & 123.456778 \\
  Texto de exemplo & abcdef ghjijklmn & 123.456778 \\
  Texto de exemplo & abcdef ghjijklmn & 123.456778 \\
  Texto de exemplo & abcdef ghjijklmn & 123.456778 \\
  Texto de exemplo & abcdef ghjijklmn & 123.456778 \\
  Texto de exemplo & abcdef ghjijklmn & 123.456778 \\
  Texto de exemplo & abcdef ghjijklmn & 123.456778 \\
  Texto de exemplo & abcdef ghjijklmn & 123.456778 \\
  Texto de exemplo & abcdef ghjijklmn & 123.456778 \\
  Texto de exemplo & abcdef ghjijklmn & 123.456778 \\

\end{longtable}


\section{Inclusión de código fonte}

Se precisamos incluír fragmentos de código fonte, podemos facelo, por exemplo,
da seguinte maneira:

\begin{lstlisting}[language=C]
#include <stdio.h>
#define N 10

int main()
{
  int i;

  // Isto é un comentario
  puts("Ola, mundo!");

  for (i = 0; i < N; i++)
  {
    puts("LaTeX é a ferramenta de edición ideal para profesionais da informática!");
  }

  return 0;
}
\end{lstlisting}

\section{Uso da relación de acrónimos e do glosario}

Os acrónimos edítanse no ficheiro \texttt{bibliografia/acronimos.tex} e úsanse
empregando a orde \texttt{acrlong} para obter o termo completo (deste xeito:
\acrlong{erlang}), a orde \texttt{acrshort} para obter o acrónimo (deste xeito:
\acrshort{erlang}). A primeira vez que usamos un termo con acrónimo no
documento é recomendable usar orde \texttt{acrfull} (que produce ambas versións
á vez: \acrfull{erlang}). Os acrónimos que non se usan no documento, non
aparecen na relación que se xerar na versión PDF.

Pola súa banda, os termos do glosario edítanse no ficheiro
\texttt{bibliografia/glo\-sa\-rio.tex} e úsanse empregando a orde \texttt{gls}
(deste xeito, \gls{bytecode}) ou \texttt{Gls} (deste xeito, \Gls{bytecode}). Ao
igual que os acrónimos, os termos que non se usan no documento, non aparecen na
relación que se xera na versión PDF.

%\include{contido/...}
\chapter{Conclusións}
\label{chap:conclusions}

\lettrine{D}{erradeiro} capítulo da memoria, onde se presentará a
situación final do traballo, as leccións aprendidas, a relación coas
competencias da titulación en xeral e a mención en particular,
posibles liñas futuras,\dots

\Blindtext


%%%%%%%%%%%%%%%%%%%%%%%%%%%%%%%%%%%%%%%%
% Apéndices, glosarios e bibliografía  %
%%%%%%%%%%%%%%%%%%%%%%%%%%%%%%%%%%%%%%%%

\appendix
\appendixpage\
\chapter{Material adicional}
\label{chap:adicional}
%\include{anexos/...}

\printglossary[type=\acronymtype,title=\nomeglosarioacronimos]
\printglossary[title=\nomeglosariotermos]

\bibliographystyle{IEEEtranN}
\bibliography{\bibconfig,bibliografia/bibliografia}
\clearpage

\end{document}

%%%%%%%%%%%%%%%%%%%%%%%%%%%%%%%%%%%%%%%%%%%%%%%%%%%%%%%%%%%%%%%%%%%%%%%%%%%%%%%%
